%%%%%%%%%%%%%%%%%%%%%%%%%%%% Document Setup %%%%%%%%%%%%%%%%%%
\documentclass[10pt]{article}
\usepackage[T2A,T1]{fontenc}
\usepackage[utf8]{inputenc}
\usepackage[russian,english]{babel}
\usepackage{CJKutf8}
\AtBeginDvi{\input{zhwinfonts}}
%\usepackage{polyglossia}
\makeatletter
\newlength{\bibhang}
\setlength{\bibhang}{1em}
\newlength{\bibsep}
 {\@listi \global\bibsep\itemsep \global\advance\bibsep by\parsep}
%\newenvironment{bibsection}%
 %       {\vspace{-\baselineskip}\begin{list}{}{%
   %    \setlength{\leftmargin}{\bibhang}%
     %  \setlength{\itemindent}{-\leftmargin}%
       %\setlength{\itemsep}{\bibsep}%
       %\setlength{\parsep}{\z@}%
        %\setlength{\partopsep}{0pt}%
        %\setlength{\topsep}{0pt}}}
        %{\end{list}\vspace{-.6\baselineskip}}
\makeatother
\reversemarginpar
\usepackage{geometry}\geometry{a4paper,total={170mm,257mm},left=1.18in,right=1.18in,top=0.79in,bottom=0.60in}
%\usepackage[paper=letterpaper,
           %\includefoot, % Uncomment to put page number above margin
   %         marginparwidth=0.6in, % Length of section titles
         %   marginparsep=0.05in, % Space between titles and text
      %      margin=1.18in, % 1 inch margins
            %includemp]{geometry}
%\usepackage{eurosans}
\setlength{\parindent}{0in}
\usepackage{enumerate} 
\usepackage{etaremune}
\usepackage{calc}
\usepackage{bbding}
\usepackage{paralist}
\usepackage{fancyhdr,lastpage}
\pagestyle{fancy}
%\pagestyle{empty} % Uncomment this to get rid of page numbers
\fancyhf{}\renewcommand{\headrulewidth}{0pt}
\fancyfootoffset{\marginparsep+\marginparwidth}
\newlength{\footpageshift}
\setlength{\footpageshift}
          {0.5\textwidth+1.\marginparsep+0.5\marginparwidth-2in}
\lfoot{\hspace{\footpageshift}%
       \parbox{4in}{\, \hfill %
                    \arabic{page} of \protect\pageref*{LastPage} % +LP
% \arabic{page} % -LP
                    \hfill \,}}
\usepackage{color,hyperref}
\usepackage{graphics}
\usepackage{wrapfig}
\usepackage{longtable}
\usepackage[usenames,dvipsnames]{xcolor}
\definecolor{darkblue}{rgb}{0.0,0.0,0.3}
\definecolor{mygreen}{RGB}{44,204,148}
\definecolor{darkgreen}{RGB}{13,107,27}
\definecolor{lightgreen}{RGB}{0,170,37}
\definecolor{provinces}{RGB}{11,100,132}
\definecolor{mypurple}{RGB}{77,24,137}
\definecolor{haiyan}{RGB}{8,150,178}
\definecolor{anticred}{RGB}{165,62,11}
\hypersetup{colorlinks,breaklinks,
            linkcolor=darkblue,urlcolor=darkblue,
            anchorcolor=darkblue,citecolor=darkblue}
%%%%%%%%%%%%%%%%%%%%%%%% End Document Setup %%%%%%%%%%%%%%%%%%%
%%%%%%%%%%%%%%%%%%%%%%%%%%% Helper Commands %%%%%%%%%

\newcommand{\makeheading}[1]%
        {\hspace*{-\marginparsep minus \marginparwidth}%
         \begin{minipage}[t]{\textwidth+\marginparwidth+\marginparsep}%
                {\large \bfseries #1}\\[-0.15\baselineskip]%
                 \rule{\columnwidth}{1pt}%
         \end{minipage}}
\renewcommand{\section}[2]%
        {\pagebreak[2]\vspace{1.3\baselineskip}%
         \phantomsection\addcontentsline{toc}{section}{#1}%
         \hspace{0in}%
         \marginpar{
         \raggedright \scshape #1}#2}
% An itemize-style list with lots of space between items
\newenvironment{outerlist}[1][\enskip\textbullet]%
        {\begin{itemize}[#1]}{\end{itemize}%
         \vspace{-.6\baselineskip}}
% An environment IDENTICAL to outerlist that has better pre-list spacing
% when used as the first thing in a \section
\newenvironment{lonelist}[1][\enskip\textbullet]%
        {\vspace{-\baselineskip}\begin{list}{#1}{%
        \setlength{\partopsep}{0pt}%
        \setlength{\topsep}{0pt}}}
        {\end{list}\vspace{-.6\baselineskip}}

% An itemize-style list with little space between items
\newenvironment{innerlist}[1][\enskip\textbullet]%
        {\begin{compactitem}[#1]}{\end{compactitem}}
% An environment IDENTICAL to innerlist that has better pre-list spacing
% when used as the first thing in a \section
\newenvironment{loneinnerlist}[1][\enskip\textbullet]%
        {\vspace{-\baselineskip}\begin{compactitem}[#1]}
        {\end{compactitem}\vspace{-.6\baselineskip}}
% To add some paragraph space between lines.
% This also tells LaTeX to preferably break a page on one of these gaps
% if there is a needed pagebreak nearby.
\newcommand{\blankline}{\quad\pagebreak[2]}
% Uses hyperref to link DOI
\newcommand\doilink[1]{\href{http://dx.doi.org/#1}{#1}}
\newcommand\doi[1]{doi:\doilink{#1}}
% For \url{SOME_URL}, links SOME_URL to the url SOME_URL
\providecommand*\url[1]{\href{#1}{#1}}
% Same as above, but pretty-prints SOME_URL in teletype fixed-width font
\renewcommand*\url[1]{\href{#1}{\texttt{#1}}}

\include{tikz-3dplot}
\usepackage{tikz,tikz-3dplot}
\usepackage{amssymb}
\usepackage{xifthen}
\usepackage{pstricks,pst-solides3d}
\usetikzlibrary{chains}
\usepackage{pgfplots}
\pgfplotsset{
compat=1.16,
    compat/path replacement=1.5.1,
}
\usepackage{pgfmath}
\usetikzlibrary{matrix}
\usetikzlibrary{arrows,decorations.pathmorphing,backgrounds,positioning,fit,petri}
\usepgfmodule{decorations}
\usepgflibrary{decorations.shapes} 
\usetikzlibrary[decorations.shapes] 
\usetikzlibrary{mindmap}
\usepgflibrary{patterns}
\usetikzlibrary[patterns]
\usetikzlibrary{petri}
\usepgflibrary{plothandlers}
\usetikzlibrary{plothandlers}
\usepackage{pgfpages}
\usepgflibrary{arrows}
\usetikzlibrary{arrows}
\usetikzlibrary{trees}
\usetikzlibrary{topaths}
\usetikzlibrary{3d}
\usetikzlibrary{positioning}
\usetikzlibrary{fit}
\usepgflibrary{shapes.misc}
\usetikzlibrary{shapes.misc}
\usepgflibrary{shadings}
\usetikzlibrary{shadings}
\usepgfmodule{shapes}
\usepgfplotslibrary{external}
\usetikzlibrary{pgfplots.external}
\usepgfplotslibrary{colorbrewer}
\usepgfplotslibrary{patchplots}
\usetikzlibrary{plotmarks}
\usepackage{color} 
\usepackage{xcolor} 
\pgfplotsset{colormap={cool}{rgb255(0cm)=(255,255,255); rgb255(1cm)=(0,128,255); rgb255(2cm)=(255,0,255)}}
  \pgfplotsset{colormap={violet}{rgb255=(25,25,122) color=(white) rgb255=(238,140,238)}}
  \pgfplotsset{  colormap={bluered}{rgb255(0cm)=(0,0,180); rgb255(1cm)=(0,255,255); rgb255(2cm)=(100,255,0);rgb255(3cm)=(255,255,0); rgb255(4cm)=(255,0,0); rgb255(5cm)=(128,0,0)}}
 \pgfplotsset{/pgfplots/colormap={winter}{rgb255=(0,0,255) rgb255=(0,255,128)}}
 \pgfplotsset{/pgfplots/colormap={spring}{rgb255=(255,0,255) rgb255=(255,255,0)}}
 \pgfplotsset{colormap={CM}{rgb=(0,0,1) color=(black) rgb255=(238,140,238)}}
 
\pgfplotsset{every patch/.style={miter limit=50},}
\usepgfplotslibrary{colormaps}
\usetikzlibrary{spy}
\begin{document}
\selectlanguage{russian}

\begin{tikzpicture} [spy using outlines={red,circle,magnification=1.5, size=2.5cm},connect spies] % <-- in the picture environment we indicate that we are going to use 'spy' zooming function, and indicate its parameters (here: red color, circle size, use line connecting it with the original place of fragment)
\begin{axis}[colorbar sampled line,small,title=Experiment №5833,xlabel={\tiny{Time, h}},
    ylabel={\tiny{Immersion depth, mm}}, zlabel={\tiny{($C_{eq}$) , MPa}},xlabel style={sloped like x axis},
    ylabel style={sloped},colormap/RdYlBu-4,grid=major,view={240}{30},minor x tick num=10,minor y tick num=10]
   		 \addplot3+ [mesh,scatter,
        samples=42,domain=0:1, ] coordinates {
 (0,1.0,0.100) (20,1.60,0.080) (40,1.70,0.060) (60,1.75,0.030) (80,1.80,0.020) (100,1.81,0.010)   
 (0,0.5,0.25) (2,0.55,0.20) (6,0.57,0.18) (8,0.62,0.17) (12,0.66,0.16) (18,0.7,0.15) 
 (0,0.300,0.30) (2,0.350,0.26) (4,0.400,0.25) (6,0.440,0.24) (8,0.490,0.22) (10,0.550,0.20) 
 (0,0.31,0.32) (2,0.33	,0.27) (4,0.41,0.24) (6,0.43,0.22) (8,0.47,0.20) (10,0.53,0.18) 
 (0,0.51,0.24) (2,0.54,0.22) (4,0.57,0.21) (6,0.58,0.18) (8,0.60,0.17) (10,0.62,0.16) 
 (0,0.2,0.23) (2,0.350	,0.21) (4,0.47,0.20) (6,0.57,0.19) (8,0.58,0.18) (10,0.61,0.17)  
};
 	       	\coordinate (spypoint) at (50,0.16,0.16);% <-- here place the 3D coordinates where it should be zoomed
  		      \coordinate (magnifyglass) at (90,1.0,0.16); % <-- here indicate 3D coordinates of the point where to place zoomed fragment
	\end{axis}
			\spy on (spypoint) in node at (magnifyglass); % <-- zoom spy command
\end{tikzpicture}
	\hskip 10pt 	
\begin{tikzpicture} [spy using outlines={red,circle,magnification=1.5, size=2.5cm},connect spies]
	\begin{axis}[colorbar sampled line,small,title=Experiment №5804,xlabel={\tiny{Time, h}},
    ylabel={\tiny{Immersion depth, mm}}, zlabel={\tiny{($C_{eq}$), MPa}}, xlabel style={sloped like x axis},
    ylabel style={sloped},colormap name=hot,grid=major,view={210}{30},minor x tick num=10,minor y tick num=10] %/RdGy-9
	    \addplot3+ [mesh,scatter,
        samples=42,domain=0:1, ] coordinates {
(0,0.5,0.09) (50,1.3,0.04) (100,1.5,0.03) (150,1.6,0.025) (200,1.7,0.020) (250,1.8,0.150) (300,1.9,0.100) 
(0,0.8,0.8) (2,0.92,0.75) (3	,0.95,0.73) (4,0.98,0.65) (5,1.0,0.62) (6,1.05,0.61) (8,1.1,0.59) 
(0,0.55,0.13) (10,0.70,0.10) (20,0.75,0.09) (30,0.80,0.083) (40,0.85,0.077) (50,0.95,0.075) (70,1.00,0.07) 
(0,0.4,0.13) (2,0.5,0.12) (4	,0.55,0.17) (6,0.57,0.14) (8,0.61,0.10) (10,0.65,0.90) (12,0.70,0.083)
(0,0.4,0.140) (2,0.5,0.130) (4,0.52,0.122) (6,0.57,0.115) (8,0.63	,0.080) (10,0.68,0.09) (12,0.70,0.080) 
(0,0.4,0.130) (2,0.5,0.128) (4,0.52,0.120) (6,0.57,0.111) (8,0.63,0.079) (10,0.68,0.086) (12,0.70	,0.077) 
 };
   	     	\coordinate (spypoint) at (50,0.6,0.16);
     	   \coordinate (magnifyglass) at (280,1.0,0.5);
	\end{axis}
		\spy on (spypoint) in node at (magnifyglass);
\end{tikzpicture}
\end{document}